% !TeX root = main.tex
% !TeX spellcheck = pt_BR


% ============================================================
% pacotes_notacao.tex
% Pacotes, notação e ambientes consolidados
% (originalmente espalhados em main.tex, inferencia.tex e prob_avancada.tex)
% ============================================================

% ---------- Pacotes ----------
\usepackage[utf8]{inputenc} 			% codificação de entrada
\usepackage[T1]{fontenc} 				% codificação de saída
\usepackage[brazil]{babel}
\usepackage{hyperref}
\usepackage{mathtools} 				% já carrega amsmath
\usepackage{amsthm}
\usepackage{amssymb}
\usepackage[brazilian]{cleveref}		% citações inteligentes de fig., eq., etc.
\usepackage{graphicx}
\usepackage{cancel} 					% cancelamento de variáveis em fórmulas
\usepackage{xcolor}					% cores em textos
\usepackage{outlines}					% itemize melhorado
\usepackage{etoolbox} 					% ajuste de espaçamento de ambientes
\usepackage{needspace}
\usepackage{bbm} 						% \mathbbm{1} (função indicadora)
\usepackage{bm}							% negrito para vetores/matrizes
\usepackage{accents}					% \prior, \post

\numberwithin{equation}{section} 		% Equações no formato X.Y

% ---------- Notação de vetores, variáveis aleatórias e matrizes ----------
\newcommand{\vect}[1]{\bm{#1}}  	    	% Vetor não aleatório
\newcommand{\matr}[1]{\bm{\mathsf{#1}}} 	% Matriz
\newcommand{\rv}[1]{#1}             		% Variável aleatória (itálico padrão)
\newcommand{\rvect}[1]{\boldsymbol{#1}} 	% Vetor aleatório
\newcommand{\prior}[1]{\accentset{\smile}{#1}} % quantidade a priori
\newcommand{\post}[1]{\accentset{\frown}{#1}}  % quantidade a posteriori

% ---------- Símbolos e conjuntos ----------
\newcommand{\cp}[1]{#1^\mathsf{c}}			% complementar de um conjunto (A^\mathsf{c})
\newcommand{\inv}[1]{#1^{-1}}      		% função ou matriz inversa
\renewcommand{\Re}{\mathbb{R}}          	% conjunto dos Reais
\newcommand{\Ree}{\widebar{\mathbb{R}}} 	% conjunto dos Reais Estendidos
\renewcommand{\P}{\mathbb{P}} 				% medida probabilidade
\newcommand{\Q}{\mathbb{Q}} 				% conjunto dos Racionais
\newcommand{\N}{\mathbb{N}}				% conjunto dos Naturais
\newcommand{\Bo}{\mathcal{B}(\Re)} 				% sigma-álgebra de Borel em \Re
\newcommand{\Born}{\mathcal{B}(\Re^n)} 			% sigma-álgebra de Borel em \Re^n
\newcommand{\Boe}{\mathcal{B}(\widebar{\Re})}		% sigma-álgebra de Borel em \Re estendido
\newcommand{\indep}{\perp\!\!\!\!\perp}			% independência

% atalhos para letras caligráficas
\newcommand{\Acal}{\mathcal{A}}
\newcommand{\Bcal}{\mathcal{B}}
\newcommand{\Ccal}{\mathcal{C}}
\newcommand{\Dcal}{\mathcal{D}}
\newcommand{\Ecal}{\mathcal{E}}
\newcommand{\Fcal}{\mathcal{F}}
\newcommand{\Gcal}{\mathcal{G}}
\newcommand{\Hcal}{\mathcal{H}}
\newcommand{\Ical}{\mathcal{I}}
\newcommand{\Mcal}{\mathcal{M}}
\newcommand{\Pcal}{\mathcal{P}}
\newcommand{\Qcal}{\mathcal{Q}}
\newcommand{\Rcal}{\mathcal{R}}
\newcommand{\Scal}{\mathcal{S}}
\newcommand{\Xcal}{\mathcal{X}}
\newcommand{\Ycal}{\mathcal{Y}}
\newcommand{\sig}{$\sigma$}

% ---------- Distribuições ----------
\DeclareMathOperator{\bernop}{Bernoulli}
\newcommand{\berndist}[1]{\bernop(#1)}

\DeclareMathOperator{\betaop}{Beta}
\newcommand{\betadist}[1]{\betaop(#1)}

\DeclareMathOperator{\binop}{Binomial}
\newcommand{\bindist}[1]{\binop\left(#1\right)}

\DeclareMathOperator{\negbinop}{Binomial-Negativa}
\newcommand{\negbindist}[1]{\negbinop(#1)}

\DeclareMathOperator{\dirichletop}{Dirichlet}
\newcommand{\dirichletdist}[1]{\dirichletop(#1)}

\DeclareMathOperator{\expop}{Exp}
\newcommand{\expdist}[1]{\expop(#1)}

\DeclareMathOperator{\gammaop}{Gamma}
\newcommand{\gammadist}[1]{\gammaop(#1)}

\DeclareMathOperator{\invgammaop}{Gamma-Inversa}
\newcommand{\invgammadist}[1]{\invgammaop(#1)}

\DeclareMathOperator{\normalgammaop}{Normal-Gamma}
\newcommand{\ngdist}[1]{\normalgammaop(#1)}

\DeclareMathOperator{\normalop}{\mathcal{N}}
\newcommand{\normdist}[1]{\normalop\left(#1\right)}

\DeclareMathOperator{\multop}{Multinomial}
\newcommand{\multdist}[1]{\multop\left(#1\right)}

\DeclareMathOperator{\poiop}{Poisson}
\newcommand{\poidist}[1]{\poiop(#1)}

\DeclareMathOperator{\tsudentop}{t-Student}
\newcommand{\tdist}[1]{\tsudentop(#1)}

\DeclareMathOperator{\triangop}{Triangular}
\newcommand{\triangdist}[1]{\triangop(#1)}

\DeclareMathOperator{\unifop}{Uniforme}
\newcommand{\unifdist}[1]{\unifop(#1)}

% ---------- Convergência ----------
\newcommand{\toQC}{\xrightarrow{\text{q.c.}}} % quase certamente
\newcommand{\toP}{\xrightarrow{\mathcal{P}}}  % em probabilidade
\newcommand{\toL}{\xrightarrow{L^p}}          % em norma L^p
\newcommand{\toD}{\xrightarrow{\mathcal{D}}}  % em distribuição

% ---------- Variável barrada ----------
\makeatletter
\newcommand{\widebar}[1]{\mathop{\overline{#1}}\nolimits\@ifnextchar^{}{\,}}
\makeatother

% ---------- Operadores ----------
\DeclareMathOperator{\E}{\mathbb{E}}     	% esperança
\DeclareMathOperator{\Var}{Var}				% variância
\DeclareMathOperator{\Cov}{Cov} 		 	% covariância
\DeclareMathOperator{\cov}{Cov}				% covariância (minúsculo)
\DeclareMathOperator{\lik}{\mathcal{L}}  	% verossimilhança
\DeclareMathOperator{\llik}{\ell}        	% log-verossimilhança
\DeclareMathOperator{\emv}{emv}				% estimador de máxima verossimilhança
\DeclareMathOperator{\map}{map}				% máximo a posteriori
\DeclareMathOperator{\kl}{KL}				% divergência Kullback-Leibler
\DeclareMathOperator{\T}{^\intercal} 		% operador de transposto (para matrizes)
\DeclareMathOperator{\ind}{\mathbbm{1}}		% função indicadora
\DeclarePairedDelimiter\abs{\lvert}{\rvert}
\DeclarePairedDelimiter\norm{\lVert}{\rVert}
\DeclareMathOperator*{\argmin}{arg\,min}
\DeclareMathOperator*{\argmax}{arg\,max}

% Swap the definition of \abs* and \norm*, so that \abs
% and \norm resizes the size of the brackets, and the
% starred version does not.
\makeatletter
\let\oldabs\abs
\def\abs{\@ifstar{\oldabs}{\oldabs*}}
%
\let\oldnorm\norm
\def\norm{\@ifstar{\oldnorm}{\oldnorm*}}
\makeatother

% ---------- Ambientes de teorema ----------


\newtheoremstyle{thm_estilo}% name
{3pt}       % espaço acima
{3pt}       % espaço abaixo
{}          % fonte do corpo
{}          % identação
{\bfseries} % fonte do cabeçalho
{.}         % pontuação após o cabeçalho
{.5em}      % espaço após o cabeçalho
{\thmname{#1}\thmnumber{ #2}\thmnote{\ \normalfont(\textit{#3})}}

\newtheoremstyle{thm_note_est}% name
{3pt}       % espaço acima
{3pt}       % espaço abaixo
{}          % fonte do corpo
{}          % identação
{\bfseries} % fonte do cabeçalho
{:}         % pontuação após o cabeçalho
{.5em}      % espaço após o cabeçalho
{\thmname{#1}\thmnumber{ #2}\thmnote{\ \normalfont(\textit{#3})}}

\theoremstyle{thm_estilo}
\newtheorem{theorem}{Teorema}[section]
\newtheorem{definition}{Definição}[section]
\newtheorem{lemma}{Lema}[section]
\newtheorem{corollary}{Corolário}[section]

\theoremstyle{thm_note_est}
\newtheorem*{note}{Nota}

\BeforeBeginEnvironment{theorem}{\vspace{1em}}     % 1em = espaço ajustável
\BeforeBeginEnvironment{definition}{\vspace{1em}}
\BeforeBeginEnvironment{lemma}{\vspace{1em}}
\BeforeBeginEnvironment{corollary}{\vspace{1em}}

% ---------- Ajustes de listas ----------
\renewcommand{\labelenumi}{(\theenumi)}
\renewcommand{\outlineii}{enumerate}
